% chapters/requirements/rf_security.tex
% RF-16: Seguridad y auditoría
% ============================================================================

\item \textbf{Seguridad y auditoría.} \label{rf:security}
El sistema debe garantizar la protección de los datos sensibles que maneja (\acs{dni},
calificaciones, datos de posibles menores) mediante medidas técnicas específicas. El
\textit{log} de auditoría es un registro inmutable de escritura única que ningún rol puede
modificar ni eliminar, constituyendo un requisito de \textit{compliance} y seguridad
operativa desde la versión~1.0. El alcance del \textit{log} de auditoría es por
organización, con un \textit{log} global para el superadministrador del sistema.

\begin{functional}

    % RF-16.1 ---------------------------------------------------------------
    \item \textbf{Registro inmutable de eventos de auditoría.}
    \label{rf:sec:audit}
    El sistema debe registrar de forma inmutable los eventos de negocio con impacto en
    datos.

    \textbf{Entrada:} evento interno generado por cualquier operación auditada del sistema.

    \textbf{Proceso:} el sistema registra cada evento en una tabla \textit{append-only} de
    PostgreSQL con permisos restrictivos a nivel de base de datos: la aplicación solo tiene
    permiso de inserción, nunca de actualización ni borrado. Cada registro incluye:
    \texttt{organization\_id}, tipo de evento, usuario actor, dirección IP de origen, marca
    temporal, entidad afectada, identificadores de la entidad y los datos relevantes del
    evento (valores anteriores y nuevos en caso de modificación). Los eventos auditados
    incluyen: inicio y cierre de sesión (éxito y fallo), acceso a \acs{pdf} de examen,
    lectura de datos sensibles (\acs{dni} descifrado), creación, modificación y eliminación
    de entidades principales (usuarios, asignaturas, exámenes, instancias), modificación de
    calificaciones, cambios de permisos y roles, operaciones de importación y exportación
    masiva, cambio de curso académico, modificación de configuración de organización y
    ejecución de borrado automático.

    \textbf{Salida:} registro de auditoría almacenado de forma permanente e inmutable.

    % RF-16.2 ---------------------------------------------------------------
    \item \textbf{Consulta del \textit{log} de auditoría.} \label{rf:sec:audit_query}
    El sistema debe permitir a un gestor consultar el \textit{log} de auditoría de su
    organización con filtros avanzados.

    \textbf{Entrada:} petición GET al recurso del \textit{log} de auditoría con parámetros
    opcionales de filtrado: rango de fechas, tipo de evento, usuario actor, entidad
    afectada, dirección IP y paginación.

    \textbf{Proceso:} el sistema recupera los registros de la organización del gestor que
    cumplen los filtros, ordenados cronológicamente (descendente por defecto). Los registros
    del \textit{log} de auditoría nunca se eliminan ni archivan: persisten de forma
    indefinida.

    \textbf{Salida:} respuesta con la lista paginada de registros de auditoría.

    % RF-16.3 ---------------------------------------------------------------
    \item \textbf{Cifrado de \acs{dni} en reposo.} \label{rf:sec:dni}
    El sistema debe cifrar el \acs{dni} de los usuarios antes de almacenarlo en la base de
    datos.

    \textbf{Entrada:} operación de creación o modificación de un usuario que incluya el
    campo \acs{dni}.

    \textbf{Proceso:} el sistema cifra el valor del \acs{dni} mediante AES-256-GCM con una
    clave maestra almacenada como variable de entorno del sistema, nunca en la base de
    datos. Cada registro utiliza un \textit{nonce} único garantizado por el modo GCM. El
    \acs{dni} cifrado se almacena junto con el \textit{nonce} y la etiqueta de
    autenticación. El descifrado se realiza solo cuando es necesario mostrar el \acs{dni} a
    un usuario autorizado (gestor de la organización o el propio usuario). La rotación de la
    clave maestra se contempla para la versión~2.0.

    \textbf{Salida:} \acs{dni} almacenado cifrado en la base de datos.

    % RF-16.4 ---------------------------------------------------------------
    \item \textbf{\textit{Hash} seguro de contraseñas.} \label{rf:sec:hash}
    El sistema debe almacenar las contraseñas de los usuarios como \textit{hash} con
    \textit{salt}, utilizando un algoritmo resistente a ataques de fuerza bruta.

    \textbf{Entrada:} contraseña en texto plano proporcionada durante la creación,
    importación o restablecimiento de contraseña.

    \textbf{Proceso:} el sistema genera un \textit{salt} aleatorio y aplica un algoritmo de
    \textit{hashing} robusto (Argon2 o bcrypt) para producir el \textit{hash}. Solo se
    almacenan el \textit{hash} y el \textit{salt}; la contraseña en texto plano se descarta
    inmediatamente después del \textit{hashing}. Durante la autenticación, se aplica el
    mismo algoritmo a la contraseña proporcionada y se compara con el \textit{hash}
    almacenado.

    \textbf{Salida:} contraseña almacenada como \textit{hash} seguro.

    % RF-16.5 ---------------------------------------------------------------
    \item \textbf{\acsp{url} temporales con expiración para \acs{pdf}.}
    \label{rf:sec:presigned}
    El sistema debe generar \acsp{url} temporales con expiración para el acceso a los
    ficheros \acs{pdf} almacenados en MinIO.

    \textbf{Entrada:} petición de acceso a un \acs{pdf} (de instancia o de modelo).

    \textbf{Proceso:} en lugar de servir el fichero directamente a través de la \acs{api},
    el sistema genera una \acs{url} \textit{presigned} de MinIO con un tiempo de expiración
    configurable por organización. La \acs{url} caduca automáticamente tras el tiempo
    configurado. Se registra el acceso en el \textit{log} de auditoría.

    \textbf{Salida:} respuesta con la \acs{url} temporal. Tras la expiración, la \acs{url}
    deja de ser válida.

    % RF-16.6 ---------------------------------------------------------------
    \item \textbf{Marca de agua dinámica en \acs{pdf} servidos a alumnos.}
    \label{rf:sec:watermark}
    El sistema debe aplicar una marca de agua dinámica a los \acs{pdf} de instancias cuando
    se sirven a alumnos, dificultando su redistribución.

    \textbf{Entrada:} petición de acceso al \acs{pdf} de una instancia por parte de un
    alumno.

    \textbf{Proceso:} el sistema genera una versión del \acs{pdf} con una marca de agua que
    incluye datos identificativos del alumno (nombre, \acs{nia}) y una marca temporal. La
    generación de la marca de agua se realiza en tiempo de petición o mediante caché para
    evitar regeneración repetida durante la misma sesión de revisión. El \acs{pdf} original
    en MinIO no se modifica.

    \textbf{Salida:} \acs{pdf} con marca de agua servido al alumno.

    % RF-16.7 ---------------------------------------------------------------
    \item \textbf{Protección anti-descarga de \acs{pdf}.}
    \label{rf:sec:anti_download}
    El sistema debe dificultar la descarga no autorizada de \acs{pdf} por parte de los
    alumnos, aplicando múltiples capas de protección.

    \textbf{Entrada:} petición de visualización de \acs{pdf} por un alumno (cuando el
    permiso de descarga está desactivado).

    \textbf{Proceso:} el sistema aplica las siguientes medidas combinadas: \acsp{url}
    temporales con expiración corta (\cref{rf:sec:presigned}), marca de agua dinámica
    (\cref{rf:sec:watermark}), cabeceras \acs{http} anti-caché y cifrado simple del
    \textit{payload} del \acs{pdf} en tránsito (cifrado con una clave derivada del
    identificador de sesión del usuario, que el \dfn{frontend} descifra en cliente). Este
    cifrado no es criptográficamente robusto contra un atacante sofisticado, pero obliga a
    inspeccionar el código del \dfn{frontend} para acceder al contenido, disuadiendo la
    descarga casual. La cabecera \texttt{Content-Disposition} se establece como
    \texttt{inline} (nunca \texttt{attachment}).

    \textbf{Salida:} \acs{pdf} servido con todas las protecciones activas.

\end{functional}
